\documentclass[12pt]{article}
\usepackage[T1]{fontenc}
\usepackage[utf8]{inputenc}
\usepackage{lmodern}
\usepackage{textcomp}
\usepackage{enumitem}
\usepackage{geometry}
\geometry{margin=1in}
\begin{document}
\begin{center}
Answer Key - Chemistry - Graduate Level Assessment 15 \
\small (Answers are ordered exactly as in the question paper.)
\end{center}

\section*{Multiple Choice}
\begin{enumerate}[label=\arabic*.]
\item \textbf{Answer:} A
\\ \textit{Options:} A) RbBr only; B) RbBr, Rb, Br$_{2}$, and Rb$_{2}$Br; C) RbBr and Rb$_{2}$Br only; D) RbBr, Rb, and Br$_{2}$; E) Rb and Br$_{2}$ only
\\ \textit{Note:} The correct answer is RbBr only because the reaction between rubidium (Rb) and bromine (Br) produces rubidium bromide (RbBr) as the sole product, and since both reactants are present in equal masses, the limiting reactant is not applicable, ensuring that only RbBr remains after the reaction completes.
\item \textbf{Answer:} D
\\ \textit{Options:} A) The hydration energies of cerium ions and sulfate ions are very high.; B) The heat of solution of cerium(III) sulfate is isothermic.; C) The solubility of cerium(III) sulfate increases with temperature.; D) The heat of solution of cerium(III) sulfate is exothermic.; E) The hydration energies of cerium ions and sulfate ions are very low.
\\ \textit{Note:} The correct answer is based on the principle that if a solute is less soluble in hot water than in cold, it indicates that the dissolution process releases heat (exothermic), thus favoring solubility at lower temperatures.
\item \textbf{Answer:} A
\\ \textit{Options:} A) 4 $10^{39}$; B) 6 $10^{36}$; C) 7 $10^{37}$; D) 5 $10^{39}$; E) 9 $10^{41}$
\\ \textit{Note:} The ratio of the electrical force to the gravitational force between a proton and an electron is approximately \( 4 \times 10^{39} \) due to the much stronger nature of the electromagnetic force compared to the gravitational force, as described by Coulomb's law and Newton's law of gravitation.
\item \textbf{Answer:} B
\\ \textit{Options:} A) 8.67 g; B) 3.82 g; C) 5.11 g; D) 3.00 g; E) 4.45 g
\\ \textit{Note:} The correct answer of 3.82 g is derived from the stoichiometry of the Wurtz reaction, where 10.0 g of n-propyl chloride (C$_{3}$H$_{7}$Cl) reacts to form hexane (C$_{6}$H$_{14}$) with a theoretical yield of 5.46 g, and after accounting for the 70\% yield, the actual mass produced is calculated as 5.46 g x 0.70 = 3.82 g.
\item \textbf{Answer:} C
\\ \textit{Options:} A) 26.3 g mol-1; B) 500 g mol-1; C) 212 g mol-1; D) 0.122 g mol-1; E) 1.22 g mol-1
\\ \textit{Note:} The molar mass of the monoprotic weak acid is calculated by determining the moles of KOH used in the neutralization, which is found by multiplying the volume of KOH (in liters) by its molarity, and then using the ratio of moles of acid to moles of KOH (1:1 for a monoprotic acid) to find the molar mass from the mass of the acid.
\end{enumerate}

\section*{True / False}
\begin{enumerate}[label=\arabic*.]
\setcounter{enumi}{5}
\item \textbf{Answer:} True
\\ \textit{Options:} A) True; B) False
\\ \textit{Note:} The van der Waals parameter a, representing the attractive forces between molecules, is derived from the ideal gas law and is correctly expressed in SI base units as kg $m^5$ s^\{-2\} mol^\{-2\}, which corresponds to the units of pressure multiplied by volume per mole squared.
\item \textbf{Answer:} False
\\ \textit{Options:} A) True; B) False
\\ \textit{Note:} The statement is false because the coverage of paint is typically calculated based on a standard thickness of around 30 microns, and 100 microns would require significantly more paint to achieve adequate coverage.
\item \textbf{Answer:} False
\\ \textit{Options:} A) True; B) False
\\ \textit{Note:} The statement is false because the maximum water vapor capacity at 27°C is approximately 25 grams per cubic meter, resulting in a total mass of around 10 kg for the entire 400 m³ room at 60\% relative humidity, not 8.4 kg.
\item \textbf{Answer:} False
\\ \textit{Options:} A) True; B) False
\\ \textit{Note:} The statement is false because the thermal stability of the hydrides of group-14 elements actually decreases down the group, with PbH 4 being the least stable due to the increasing influence of the heavier, more electropositive lead atom, which results in weaker H-X bonds.
\item \textbf{Answer:} True
\\ \textit{Options:} A) True; B) False
\\ \textit{Note:} Sulfur dioxide (SO 2) is classified as an acid anhydride because it reacts with water to produce sulfurous acid (H$_{2}$SO$_{3}$), demonstrating its role as a precursor to an acid.
\end{enumerate}

\section*{Long Form Response}
\begin{enumerate}[label=\arabic*.]
\setcounter{enumi}{10}
\item \textbf{Answer:} To calculate the frequency of a light wave corresponding to the Bohr radius of 0.5292 \ensuremath{\times} 10^-8 cm, we first convert the wavelength into meters, yielding 5.292 \ensuremath{\times} 10^-10 m. Using the equation \(c = \lambda \nu\), where \(c\) is the speed of light (approximately 3.00 \ensuremath{\times} $10^8$ m/s), we can find the frequency (\(\nu\)) by rearranging the formula to \(\nu = c / \lambda\). Substituting the values gives us \(\nu \approx 5.665 \times 10^{17} \, \text{sec}^{-1}\). This frequency is significant in quantum mechanics as it corresponds to the energy transitions of electrons within atoms, particularly in hydrogen-like systems, illustrating the quantum nature of light and matter interactions.
\\ \textit{Note:} correct because it accurately applies the wave equation \(c = \lambda \nu\) to convert the given wavelength into frequency, demonstrating the relationship between light properties and their quantum mechanical implications, particularly in relation to atomic structure as defined by the Bohr model.
\item \textbf{Answer:} In quantum mechanics, Hermitian operators play a crucial role as they correspond to observable quantities and guarantee real eigenvalues, which are essential for physical measurements. To determine whether differential operators are Hermitian, one must check if they satisfy the condition \( \langle \psi | \hat{A} \phi \rangle = \langle \hat{A} \psi | \phi \rangle \) for all suitable wave functions \( \psi \) and \( \phi \). Operators b and c qualify as Hermitian because they fulfill this condition under appropriate boundary conditions, leading to real eigenvalues. In contrast, if an operator does not meet this criterion, it is deemed non-Hermitian, which may imply complex eigenvalues and a lack of physical observability. Thus, the analysis confirms that operators b and c are indeed Hermitian.
\\ \textit{Note:} The answer correctly identifies that Hermitian operators correspond to physical observables in quantum mechanics and must satisfy the inner product condition for all suitable wave functions, thereby validating the classification of operators b and c as Hermitian based on their adherence to this criterion.
\end{enumerate}

\vfill
\noindent\textit{End of Answer Key}
\end{document}